\section{End-to-end: how \Flux works, today and at target}
\label{sec:end-to-end}

Every preceding section has described one layer in isolation, at the depth that layer deserves. This
section does the opposite: it takes a single application through its whole life --- constructed,
paid for, validated, placed, attested, made reachable, run, renewed, possibly reported, and
ultimately the reason somebody gets paid --- and follows it across every layer at once. It does that
twice. The first pass is the network as deployed on 2026-09-03, written in the present tense
and tagged \shipped throughout. The second pass is the target architecture specified in Parts~II
through~IV, written in the conditional and tagged \planned throughout, with the same twelve steps in
the same order so that the two can be read side by side and differenced step by step. The two passes
never mix within a paragraph, and where the deployed answer is thin --- there is no fee pool, there
is no attestation a tenant can check, there is no grace period, there is no quorum --- this section
says so at the step, not in a footnote at the end.

Nothing here is derived. Every constant, mechanism and measurement is cited to the section that
established it, and where two sections disagree about a constant this section says which two and
marks the disagreement rather than picking one.

\begin{figure}[H]
\centering
\begin{tikzpicture}[
  font=\tiny, x=1mm, y=1mm,
  cell/.style={draw, rounded corners=1.5pt, text width=48mm, align=left, inner sep=2pt,
               minimum height=8mm, anchor=west},
  tcell/.style={cell, fill=black!5},
  dead/.style={cell, draw=black!45, dashed, text=black!55, align=center},
  stepno/.style={font=\tiny\bfseries, anchor=east},
  stepname/.style={font=\scriptsize, anchor=east, text width=22mm, align=right},
  hdr/.style={font=\scriptsize\bfseries, anchor=west},
]
  \node[hdr] at (5,9)  {deployed today};
  \node[hdr] at (62,9) {at target};
  \draw[black!30] (-32,4.5) -- (112,4.5);

  \node[stepno] at (-30,0) {1};
  \node[stepname] at (-26,0) {construct};
  \node[cell] (d1) at (5,0) {\specv{8} spec via \texttt{/\allowbreak{}apps/\allowbreak{}appregister}};
  \node[tcell] (t1) at (62,0) {\specv{9} spec; all applications migrated forward};
  \node[stepno] at (-30,-11) {2};
  \node[stepname] at (-26,-11) {pay};
  \node[cell] (d2) at (5,-11) {one transparent tx to \texttt{t3Nryf\ldots}, \texttt{OP\_RETURN} message hash};
  \node[tcell] (t2) at (62,-11) {one output, full price, to $\mathsf{sink}$ + \texttt{OP\_RETURN} $H(m_\app)$};
  \node[stepno] at (-30,-22) {3};
  \node[stepname] at (-26,-22) {validate};
  \node[cell] (d3) at (5,-22) {\texttt{processInsight}; underpayment silently dropped};
  \node[tcell] (t3) at (62,-22) {PS1--PS4; \textbf{consensus} applies the 80/20};
  \node[stepno] at (-30,-33) {4};
  \node[stepname] at (-26,-33) {fee pool};
  \node[dead] (d4) at (5,-33) {---};
  \node[tcell] (t4) at (62,-33) {$\pool$ committed in the coinbase; drip $D_\hgt$};
  \node[stepno] at (-30,-44) {5};
  \node[stepname] at (-26,-44) {select payees};
  \node[cell] (d5) at (5,-44) {\PoN{} rotation, one node per tier};
  \node[tcell] (t5) at (62,-44) {\PoN{} rotation, unchanged};
  \node[stepno] at (-30,-55) {6};
  \node[stepname] at (-26,-55) {place};
  \node[cell] (d6) at (5,-55) {broadcast, \SI{90}{\second} wait, earliest wins};
  \node[tcell] (t6) at (62,-55) {\specv{9} reconciler and quorum plane};
  \node[stepno] at (-30,-66) {7};
  \node[stepname] at (-26,-66) {attest};
  \node[dead] (d7) at (5,-66) {no tenant-reachable path};
  \node[tcell] (t7) at (62,-66) {bonded $k$-of-$n$ committee, machine-bound key};
  \node[stepno] at (-30,-77) {8};
  \node[stepname] at (-26,-77) {make reachable};
  \node[cell] (d8) at (5,-77) {\FDM{} / \PDNS{} / HAProxy};
  \node[tcell] (t8) at (62,-77) {unchanged; \VPON{} for private deployments};
  \node[stepno] at (-30,-88) {9};
  \node[stepname] at (-26,-88) {run, monitor};
  \node[cell] (d9) at (5,-88) {reconciler, crash backoff};
  \node[tcell] (t9) at (62,-88) {unchanged};
  \node[stepno] at (-30,-99) {10};
  \node[stepname] at (-26,-99) {renew or lapse};
  \node[cell] (d10) at (5,-99) {expiry is arithmetic};
  \node[tcell] (t10) at (62,-99) {the payment \emph{is} the renewal};
  \node[stepno] at (-30,-110) {11};
  \node[stepname] at (-26,-110) {report, evict};
  \node[cell] (d11) at (5,-110) {GitHub-file blocklist, $\kappa=1$};
  \node[tcell] (t11) at (62,-110) {on-chain report, \SI{24}{\hour}, fee-less votes, ban memo};
  \node[stepno] at (-30,-121) {12};
  \node[stepname] at (-26,-121) {compensate};
  \node[cell] (d12) at (5,-121) {issuance only};
  \node[tcell] (t12) at (62,-121) {issuance $+$ drip $D_{\hgt,t}$};

  % the one cross-track coupling
  \draw[-{Latex[length=1.6mm]}, thick, black!70]
    (t7.east) .. controls +(14mm,0) and +(14mm,0) ..
    node[font=\tiny, right, align=left, pos=0.5, xshift=1mm]
      {the eligibility gate\\is what compels hosting} (t12.east);

  \node[font=\tiny, anchor=west, text width=52mm] at (-32,-137)
    {Dashed cells mark steps with no deployed counterpart. Every other target-side cell is
     \planned; only the shaded column changes.};
\end{tikzpicture}
\caption{The twelve-step application lifecycle, deployed and at target, on a shared step axis. Two steps have no deployed counterpart at all --- there is no fee pool, and no attestation a tenant can reach. The single cross-track edge is \Cref{prop:e2e-gate}: because \PNR{} pays every node regardless of what it hosts, the eligibility gate of step~7 is the only thing that makes step~12 compensate hosting rather than mere presence.}
\label{fig:e2e-twelve-steps}
\end{figure}

% =====================================================================
\subsection{Pass one: the deployed system}
\label{sec:e2e:deployed}

Everything in this subsection is \shipped and is written in the present tense.

\paragraph{1. A tenant constructs a specification.}
The document is a small owner-signed JSON object, and the version it may declare is fixed by chain
height: \specv{8} is enforced from mainnet height \num{1932380}, and above that height a
specification whose top-level keys are anything other than
\texttt{version,\allowbreak{} name,\allowbreak{} description,\allowbreak{} owner,\allowbreak{} compose,\allowbreak{} instances,\allowbreak{} contacts,\allowbreak{} geolocation,\allowbreak{} expire,\allowbreak{} nodes,\allowbreak{}
staticip,\allowbreak{} datacenter,\allowbreak{} enterprise} is rejected outright (\S10)
\srccode{flux/\allowbreak{}ZelBack/\allowbreak{}config/\allowbreak{}default.js}{231--240}
\srccode{flux/\allowbreak{}ZelBack/\allowbreak{}src/\allowbreak{}services/\allowbreak{}appRequirements/\allowbreak{}appValidator.js}{1063--1067}. Seven versions are simultaneously live on the
network --- \specv{2} through \specv{8}, with \num{322} of \num{1195} deployed applications
(\SI{26.9}{\percent}) below \specv{8} --- because the acceptance predicate has a ceiling and no floor
(\S10) \srcmeasured{api.runonflux.io/\allowbreak{}apps/\allowbreak{}globalappsspecifications}{2026-09-03}. The tenant signs the
envelope $\{\texttt{app\allowbreak{}Specification}, \texttt{timestamp}, \texttt{signature}, \texttt{type},
\texttt{version}\}$, where \texttt{type} is \texttt{zelappregister}\slash\texttt{fluxappregister} and the
envelope version --- distinct from the specification version --- is $1$, and \texttt{POST}s it to any
node's \texttt{/\allowbreak{}apps/\allowbreak{}appregister} at \texttt{user} privilege, which is a \texttt{bitcoinMessage}
signature over a login phrase no more than \SI{16}{\hour} old (\S9)
\srccode{flux/\allowbreak{}ZelBack/\allowbreak{}src/\allowbreak{}services/\allowbreak{}appDatabase/\allowbreak{}registryManager.js}{1755--1764}. The receiving node
refuses the request if it holds fewer than $8$ outbound or $4$ inbound peers, refuses a timestamp
more than \SI{1}{\hour} old or more than \SI{5}{\minute} ahead, and refuses a \texttt{version: 9}
specification outright unless \texttt{config.\allowbreak{}fluxapps.\allowbreak{}acceptV9} is set, which defaults to
\texttt{false} with no caller anywhere that sets it (\S9, \S10). It then computes
$\texttt{messageHASH} = \mathrm{SHA\text{-}256}(\texttt{type} \| \texttt{version} \|
\mathrm{JSON}(\texttt{spec}) \| \texttt{timestamp} \| \texttt{signature})$, stores the message with a
\SI{3600}{\second} TTL, broadcasts it, blocks through a fixed confirmation round trip of
$2 \times \SI{1200}{\milli\second}$ plus up to $20 \times \SI{500}{\milli\second}$ of polling, and
returns the hash and nothing else \srccode{flux/\allowbreak{}ZelBack/\allowbreak{}src/\allowbreak{}services/\allowbreak{}appDatabase/\allowbreak{}registryManager.js}{1827--1858}.
The specification carries no \texttt{duration} field and no \texttt{paymentMemo} field; lifetime is
expressed as \texttt{expire} in blocks, and there is no field anywhere in \specv{8} that binds a
payment to a specification other than the hash the tenant is about to put in an
\texttt{OP\_RETURN}.

\paragraph{2. The tenant pays.}
Price is a pure function of the specification and the height, evaluated against a height-scheduled
table with six segments --- genesis plus five revisions --- whose current row takes effect from
height \num{1597156} at $\pricevec = (0.03,\ 0.01,\ 0.004)\flux$ per unit of
$(\mathrm{cpu}, \mathrm{ram}, \mathrm{hdd})$ with surcharges
$(p_{\mathrm{port}}, p_{\mathrm{scope}}, p_{\mathrm{sip}}) = (0.4,\ 0.8,\ 0.4)\flux$ and
$p_{\min} = 0.01\flux$ (\S16) \srcmeasured{api.runonflux.io/\allowbreak{}apps/\allowbreak{}deploymentinformation}{2026-09-03}.
The base is $B(\app,\hgt) = 10\,p_{\mathrm{cpu}}\resvec^{\mathrm{cpu}} +
\tfrac{1}{100}p_{\mathrm{ram}}\resvec^{\mathrm{ram}} + p_{\mathrm{hdd}}\resvec^{\mathrm{hdd}}$ plus
the surcharge terms, quoted per three instances --- $\beta = \lceil B/3 \rceil_2$ per instance,
scaled by the instance count from height \num{1890000} in the three-case form of \Cref{def:price}
--- then scaled by $\dur/L^{\ast}(\hgt)$ with $L^{\ast} = 4\blockslasting = \num{88000}$ blocks at
post-\PoN heights, ceiled to two decimals and only then floored at $p_{\min}$
(\S16) \srccode{flux/\allowbreak{}ZelBack/\allowbreak{}src/\allowbreak{}services/\allowbreak{}utils/\allowbreak{}appUtilities.js}{89--134}
\srccode{flux/\allowbreak{}ZelBack/\allowbreak{}src/\allowbreak{}services/\allowbreak{}appMessaging/\allowbreak{}messageVerifier.js}{736--745}. Duration is bounded to
$\dur \in \{\num{1},\dots,\num{1056000}\}$ blocks --- the floor is height-gated by the validator at
\num{5000}, then \num{100} from height \num{1630040}, then \num{1} from height \num{1964447}, while
the published endpoint still reports \num{5000}
\srccode{flux/\allowbreak{}ZelBack/\allowbreak{}src/\allowbreak{}services/\allowbreak{}appRequirements/\allowbreak{}appValidator.js}{852--862} --- and instance count to
$\ninst \in \{\num{1},\dots,\num{100}\}$ for \specv{8} applications from height \num{2176519}
(\S16). The tenant then broadcasts exactly one transparent transaction paying at least that amount to
the sink address \texttt{t3Nryf\allowbreak{}AQLGe\allowbreak{}Fs9j\allowbreak{}Eoeqsxm\allowbreak{}BN2QLRa\allowbreak{}RKFLUX}, in force from height \num{1670000},
carrying \texttt{messageHASH} in an \texttt{OP\_RETURN} (\S16)
\srccode{flux/\allowbreak{}ZelBack/\allowbreak{}config/\allowbreak{}default.js}{241--245}. All application revenue on the network lands at
that one transparent address, by policy and not by consensus (\S12, \S16).

Both apparent disagreements here are resolved. The duration denominator carries a \PoN{} fork branch --- $22{,}000$ blocks before activation, $88{,}000$ from height $2{,}020{,}000$, the multiplier preserving wall-clock duration --- and \S7's definition now states it. The instance lower bound is $3$ in general and $1$ for \specv{8} applications from height $2{,}176{,}519$; the published endpoint reports $3$ regardless, which is a defect in the endpoint rather than a disagreement between sections \srcdoc{reconciliation record}.

\paragraph{3. The payment is validated.}
No consensus rule sees any of this. Each node's own chain-scan loop,
\texttt{explorerService.\allowbreak{}processInsight}, runs per block, recognises the payment by its output
address, sums \texttt{valueSat} across matching outputs, decodes the temporary-message hash from the
\texttt{OP\_RETURN}, and queues \texttt{messageVerifier.\allowbreak{}checkAndRequestApp} (\S9, \S16)
\srccode{flux/\allowbreak{}ZelBack/\allowbreak{}src/\allowbreak{}services/\allowbreak{}explorerService.js}{309--345,444}. Promotion re-verifies update
signatures against the current permanent owner as an explicit anti-race measure, recomputes the
\PoN-fork-aware expiration height $e_\app$, and admits the specification to \texttt{globalzelapps}
if and only if $v^{\mathrm{obs}} \ge \price \cdot 10^{8}\sat$. If the payment is short, the
registration is \textbf{silently dropped}: a server-side log warning, no error surfaced to the payer,
no retry, and no refund path anywhere in the repository (\S9, \S16)
\srccode{flux/\allowbreak{}ZelBack/\allowbreak{}src/\allowbreak{}services/\allowbreak{}appMessaging/\allowbreak{}messageVerifier.js}{735--763}. The funds are spent
and the deployment does not exist. \S16 records this as the most important correctness gap on the
path this paper otherwise recommends to autonomous software, and states the signed negative
acknowledgement that would repair it as requirement R16.1.

\paragraph{4. Value enters the fee pool.}
\textbf{This step does not exist today.} There is no fee pool. There is no consensus object of any
kind mediating between what a tenant pays and what an operator receives: the payment lands at a
transparent address controlled by policy, and every satoshi any node operator earns is newly issued
(\S7). The step is listed here, empty, so that the two passes stay aligned and so that the reader
sees the gap in the same position in both.

\paragraph{5. Consensus selects the payees.}
Independently of everything above, and with no reference to it, \fluxd awards block production every
\SI{30}{\second}. The \PoN slot is $\tslot(t) = \lfloor (t - t_0)/\spacingpon \rfloor$ with
$\spacingpon = \SI{30}{\second}$, and a confirmed \FluxNode with collateral outpoint $c$ is eligible
for slot $\tslot$ iff $H(c \| h_{\mathrm{prev}} \| \tslot) \le T_\tslot$ --- sortition over the
confirmed-node set, not weighted by collateral amount beyond confirmation (\S4)
\srccode{fluxd/\allowbreak{}src/\allowbreak{}pon/\allowbreak{}pon.cpp}{70--93}. Each block pays one node from each of the three tiers,
selected from a per-tier queue ordered by \texttt{n\allowbreak{}Last\allowbreak{}Paid\allowbreak{}Height} ascending, which is a strict
longest-time-since-paid FIFO (\S4) \srccode{fluxd/\allowbreak{}src/\allowbreak{}fluxnode/\allowbreak{}fluxnode.h}{313--328}. The measured
rotation lengths are $\rot{\tC} = \num{3247}$, $\rot{\tN} = \num{1615}$ and $\rot{\tS} = \num{1627}$
blocks, i.e.\ \SI{27.06}{\hour}, \SI{13.46}{\hour} and \SI{13.56}{\hour}
\srcmeasured{api.runonflux.io/\allowbreak{}daemon/\allowbreak{}listzelnodes}{2026-09-03}. The payee set of the block that
confirms a tenant's payment has nothing to do with that tenant, that payment, or that application.

\paragraph{6. The application is placed across nodes.}
There is no scheduler. Every node runs the same loop against its own gossiped replica of
\texttt{globalzelapps} and reaches its own conclusion (\S9). A node first clears a precondition gate
--- enterprise identity resolved, daemon synced, local database ready, not DoS-flagged, node
confirmed, \fluxbench not erroring, own external IP known, and fewer than
$\texttt{maxAppsPerNode} = 200$ applications installed --- then draws a candidate from the set of
non-expired, under-replicated applications, biased first to those naming its own IP and otherwise
uniformly at random \srccode{flux/\allowbreak{}ZelBack/\allowbreak{}src/\allowbreak{}services/\allowbreak{}appLifecycle/\allowbreak{}appSpawner.js}{78--336}. After
targeting and geolocation filters and a set of soft-deferral branches, it verifies the image against
the Docker Hub whitelist, pulls it, probes port availability through its peers, broadcasts a
\texttt{fluxappinstalling} message carrying its own IP and a self-reported \texttt{broadcastedAt},
and then waits $\texttt{install\allowbreak{}Collision\allowbreak{}Wait\allowbreak{}Ms} = \SI{90}{\second}$
\srccode{flux/\allowbreak{}ZelBack/\allowbreak{}config/\allowbreak{}default.js}{321}. It re-reads the global installing and running lists,
ranks all installing nodes by \texttt{broadcastedAt} ascending, and abandons the attempt if its rank
does not fit the remaining slots. \textbf{Earliest-broadcast-wins is the entire conflict-resolution
mechanism}: no leader election, no quorum, no on-chain arbitration
\srccode{flux/\allowbreak{}ZelBack/\allowbreak{}src/\allowbreak{}services/\allowbreak{}appLifecycle/\allowbreak{}appSpawner.js}{686--703}. Installation proceeds
through \texttt{register\allowbreak{}App\allowbreak{}Locally} into Docker, and from that point every start and stop of the
container goes through one component, the reconciler, described in its own source as the single
level-based actuator, with a crash-backoff ladder of $[0,\ \num{30000},\ \num{300000},\
\num{900000},\ \num{1800000}]\,\si{\milli\second}$ reset after \SI{600000}{\milli\second} of stable
running (\S9) \srccode{flux/\allowbreak{}ZelBack/\allowbreak{}config/\allowbreak{}default.js}{130--131}. Sixty seconds after installing, the
node re-ranks itself by \texttt{runningSince} and uninstalls immediately if it is now surplus.
Convergence is not agreement; it is two opposing correction loops, and both of them correct
over-installation while neither corrects under-installation faster than the next node's random draw
(\S9).

\paragraph{7. The host is attested.}
Nothing a tenant can check. \fluxsas, the attestation service, serves a single listener on
\texttt{127.0.0.1:\allowbreak{}16100} with mutual TLS mandatory and authorisation by client-certificate common
name --- \fluxbench for every route, \texttt{fdm-nginx} for three --- so an external tenant has no
network path to it at all (\S14) \srccode{fluxsas/\allowbreak{}src/\allowbreak{}api\_server.h}{1042--1053,220--252}. Even with
a path, the signature would not say what a tenant needs: the Ed25519 key used by \texttt{/v2/attest}
for the default, \texttt{cdn}, \texttt{mesh} and \texttt{app} purposes is derived from static values
compiled identically into every \fluxsas binary, so a valid attestation demonstrates possession of
something every copy of the software contains rather than that a particular machine booted a
particular image (\S14) \srcprivate{\SAS{} API server}. One layer down, the eligibility
gate that admits a node to the paid set at all is a benchmark signature produced by a process on the
operator's own host and verified against \texttt{vec\allowbreak{}Benchmarking\allowbreak{}Public\allowbreak{}Keys}, five network-wide keys
with timestamp validity windows --- collateral binds the tier, and the benchmark signature does not
bind the hardware to it (\S4) \srccode{fluxd/\allowbreak{}src/\allowbreak{}chainparams.cpp}{233--240}. What \emph{does} do real
work is the boot chain: a Flux-controlled Secure Boot Platform Key pinned by hash inside \texttt{fes},
fail-closed exit codes \num{50}--\num{57} that power the machine off rather than continue, a TPM2
seal against PCR~7, and a dm-verity root filesystem (\S14). But the allow-list of known-good verity
root hashes is served by two Flux-operated web endpoints that \textbf{fail open} --- if both are
unreachable, the watchdog treats the node as secure and does nothing
\srcprivate{\SAS{} watchdog}. The honest summary of this step is that measured boot
imposes real cost on a remote attacker, and that the assurance a node offers is a Flux-operated one
rather than a hardware-rooted, tenant-verifiable one.

\paragraph{8. \FDM and \PDNS make the application reachable.}
The reverse path from a hostname to a container is Foundation-operated end to end (\S12). \FDM does
not talk to \fluxd or to any peer-to-peer layer: it polls one Foundation-operated REST API,
\texttt{api.\allowbreak{}runonflux.\allowbreak{}io}, for \texttt{apps/\allowbreak{}globalappsspecifications} every \SI{30}{\second} under an
\texttt{ETag}, \texttt{apps/locations} every \SI{10}{\second} unconditionally, and permanent messages
every \SI{120}{\second}, and no fallback source exists anywhere in the codebase
\srccode{flux-domain-manager/\allowbreak{}src/\allowbreak{}services/\allowbreak{}flux/\allowbreak{}dataFetcher.js}{47--91}. Names are computed by string
concatenation --- \texttt{$\app$\_$P$.app2.runonflux.io} per port and \texttt{$\app$.app2.runonflux.io}
as the main alias --- and served either by Cloudflare in strict DNS-only mode or by the self-hosted
\PDNS zone service. The \FDM host's own HAProxy terminates TLS against every \texttt{.pem} under
\texttt{/\allowbreak{}etc/\allowbreak{}ssl/\allowbreak{}fluxapps/\allowbreak{}} and routes on the \texttt{Host} header to one of the application's
container IPs, with per-server health checks of \texttt{inter 3s fall 2 rise 2 fastinter 500},
i.e.\ an approximately \SI{6}{\second} detection floor
\srccode{flux-domain-manager/\allowbreak{}src/\allowbreak{}services/\allowbreak{}haproxyTemplate.js}{260--319}. Sixteen Foundation-run
\FDM/CDM hosts serve this path for the whole network; roughly two dozen game-server application types
bypass HAProxy entirely through a single, unreplicated \texttt{flux-\allowbreak{}apps-\allowbreak{}dns-\allowbreak{}manager} instance
(\S12). Inside the fleet, container-to-container names resolve through \fluxdnsd, which despite the
name is a purely local file-tail resolver over a \texttt{membership.\allowbreak{}json} that \FluxOS writes, with a
fixed \SI{5}{\second} answer TTL and \texttt{REFUSED} rather than \texttt{SERVFAIL} when it holds no
snapshot (\S12).

\paragraph{9. It runs and is monitored.}
The reconciler holds the container to its desired state, resolved from four inputs in strict priority
order with the operator's own lock winning over all (\S9). Around it run the periodic monitors: a
node-status monitor every \SI{1.5}{\minute}, an application-availability check every
\SI{3}{\minute}, a CPU-usage check every \SI{900000}{\milli\second}, an image-compliance sweep every
\SI{3600000}{\milli\second}, and a storage check every \SI{20}{\minute} (\S9)
\srccode{flux/\allowbreak{}ZelBack/\allowbreak{}src/\allowbreak{}services/\allowbreak{}serviceManager.js}{433--563}. Location records carry a
\SI{7500}{\second} TTL, and a departing node's \texttt{fluxnodesigterm} broadcast extends its
applications' location TTL by $\texttt{SIGTERM\_EXPIRY\_MS} = \SI{420000}{\milli\second}$ rather than
expiring them at once. Where the operator has opted an application in, \texttt{flux-\allowbreak{}telemetryd} ships
per-container telemetry to a customer-chosen backend every \SI{15}{\second}. Crucially, none of this
is metering: \Flux prices reservation, not consumption, so there is no usage counter to read, to
attest, or to dispute (\S16). The only disputable fact is the binary one --- did the application run
--- and that is observable by third parties from location records, \FDM health checks, and the
tenant's own probes.

\paragraph{10. It renews or lapses.}
Expiry is arithmetic, not a decision. Entitlement is the predicate
$\mathrm{Ent}(\app,\hgt) = \indic{\hgt \le e_\app}$, where $e_\app$ was fixed at promotion from
$(\hgt_0, \dur)$ and is a function of the confirmed chain prefix alone (\S16). Every node evaluates
it independently and reaches the same answer, and a periodic pass every
$\texttt{expire\allowbreak{}Flux\allowbreak{}Apps\allowbreak{}Period} = 100$ blocks removes what has passed
\srccode{flux/\allowbreak{}ZelBack/\allowbreak{}config/\allowbreak{}default.js}{298}. A renewal is an update, priced at
$\max\{p_{\min},\ 0.9\,(\price - \varpi\price_{\mathrm{prev}})\}$ where $\varpi$ is the unused fraction
of the previous term (\S16). \textbf{There is no grace window, no suspension state, and no
data-retention guarantee}: when the paid-through height arrives the entitlement stops, and the
removal path is the ordinary one. That is a consequence of the design rather than an omission ---
Property ``No standing obligation'' in \S16 shows that any grace state would have to be a function of
information not in the chain prefix --- but it is a real difference from what a tenant used to
hyperscaler billing expects. The only automatic renewal that exists is
\texttt{appsmonitor}, a Foundation-operated daemon that renews applications whose owner keys the
Foundation itself holds, alongside a hardcoded single-address \texttt{usersToExtend} allow-list
permitted to submit expire-only extensions on any owner's behalf (\S16).

\paragraph{11. It may be reported and evicted.}
There is no quorum. Moderation is a publisher: \FluxOS fetches
\texttt{helpers/\allowbreak{}blockedrepositories.\allowbreak{}json} from \texttt{raw.\allowbreak{}githubusercontent.\allowbreak{}com/\allowbreak{}RunOnFlux/\allowbreak{}flux},
caches it, and every \SI{1}{\hour} removes every locally installed application whose image, owner
address or specification hash matches, spacing removals \SI{3}{\minute} apart (\S9, \S23)
\srccode{flux/\allowbreak{}ZelBack/\allowbreak{}src/\allowbreak{}services/\allowbreak{}appSecurity/\allowbreak{}imageManager.js}{179--195,644--676}. The wider
\texttt{fluxos-\allowbreak{}network-\allowbreak{}policy} channel carries the same shape for blocked and vetted repositories,
the enterprise-node owner map, a node-tampering blocklist and a $\sim$2.0-million-row geolocation
table, fetched every \SI{6}{\hour} or \SI{12}{\hour} and enforced fleet-wide with no canary and no
staged rollout (\S12). The documents are unsigned; by the component's own account their authenticity
rests on GitHub and on who can merge, and branch protection deliberately imposes no second-reviewer
requirement. \S23 states the resulting measure exactly: the number of parties whose agreement removes
an arbitrary application from every node is $\kappa_{\mathrm{deployed}} = 1$, the authentication on
the decision is nothing, the record a node can verify is nothing, and the latency to fleet-wide
effect is bounded by \SI{7}{\hour} (\SI{13}{\hour} for the tampering list).

\paragraph{12. Operators are compensated.}
From issuance alone. For $\hgt \ge \hpon = \num{2020000}$ the block subsidy is
$\subsidypon(\hgt) = s_{r(\hgt)}$ with $s_0 = 14\flux$, $s_{k+1} = \lfloor s_k \cdot 9/10 \rfloor$,
and $r(\hgt) = \min(\lfloor (\hgt - \hpon)/\redint \rfloor,\ \redmax)$ for
$\redint = \num{1051200}$ blocks and $\redmax = 20$ (\S5)
\srccode{fluxd/\allowbreak{}src/\allowbreak{}main.cpp}{2799--2816}. Of that, the three tier payees receive
$\lfloor \subsidypon(\hgt)\,\tbase{X}/14\flux \rfloor$ with
$(\tbase{\tC},\tbase{\tN},\tbase{\tS}) = (1,\ 3.5,\ 9)\flux$, and the remainder
$\devfund(\hgt) = 0.5\flux$ at $r = 0$ --- $1/28$, or \SI{3.57}{\percent}, of \PoN emission --- is a
consensus-enforced payment to the hardcoded address \texttt{t3h\allowbreak{}Pu1YDe\allowbreak{}GUCp8m7BQCnn\allowbreak{}NUm\allowbreak{}RMJBa5Rady\allowbreak{}A}
(\S5) \srccode{fluxd/\allowbreak{}src/\allowbreak{}main.cpp}{2877--2895}. Emission does not stop: after twenty reductions the
subsidy freezes at $\num{1.70207313}\flux$ per block, i.e.\ $\num{1789219.274256}\flux$ per year,
indefinitely (\S5). \textbf{Nothing the tenant paid in step~2 reaches any of these payees.} The
tenant's money is at one address and the operator's money is newly minted, and the only thing
connecting them is that both happen on the same chain.

% =====================================================================
\subsection{Pass two: the target architecture}
\label{sec:e2e:target}

Everything in this subsection is \planned. Nothing in it is deployed, no part of it is scheduled to a
height except where a height is explicitly named as proposed, and no sentence below should be read as
a statement about the network as it runs. Where the target depends on something unbuilt or undecided
that dependency is named at the step, in the same paragraph.

\paragraph{1. A tenant would construct a specification.}
The document would be a \specv{9} specification carrying \texttt{duration} as a first-class field, a
\texttt{paymentMemo} of up to \num{512} characters in the draft schema, a \texttt{referralCode}, and
a \texttt{machineType} drawn from $\{\texttt{standard}, \texttt{gpu}, \texttt{vps}\}$, with the four
\specv{8} renames applied and \texttt{enterprise} removed (\S10). Applications already deployed would
not be left behind on older versions and would not be offered the choice: the settled migration is
\textbf{forced}, an application-update message carrying a special signature that rewrites every
application forward onto \specv{9}, after which updates on older specification versions would be
disabled and, once the fleet had moved, the legacy validation paths removed (\S10)
\srcintent{design intake}. That supersedes the version-floor-plus-deprecation
design and is the better answer to the \num{322}-application, \SI{26.9}{\percent} legacy tail of
step~1 of the first pass, because it would clear the tail at one moment rather than letting it trail
off at each owner's own pace. It also creates a capability this paper names rather than lets a reader
find: a key able to rewrite a tenant's signed specification without that tenant's signature, which
\S10 places in the same class as the unsigned network policy, the release channel and the
emergency-block keys, and which the team's decision bounds as single-use, migration-scoped and
hardcoded --- with no opt-out for an owner who would rather remain on the document they signed
(\S10, C40).

Five things are undecided or unbuilt at this step and each is load-bearing. \specv{9} has no
enforcement height in shipped configuration or on the live pricing endpoint --- the height itself is
now settled at $\hstar = \num{3050000}$ (\S10) but is not yet shipped --- and \S10 states the
ordering constraint that the height must \emph{follow} implementation of the three economic fields
rather than precede it. \texttt{paymentMemo}, \texttt{referralCode} and \texttt{machineType} have
\textbf{zero consumer code anywhere in \FluxOS} (\S10) --- they are exactly the fields \PNR,
referrals and the VPS tier are each stated to depend on. The RFC-0 draft has no signature field at
all and states its canonicalization rules as prose \texttt{MUST}s with no schema or code enforcement,
which for a document meant to be compared byte-for-byte across roughly \num{6500} independent
validators is the defect to name (\S10). The removal of \texttt{enterprise} arrives with no declared
replacement, so \specv{9} as drafted cannot express a private deployment at all, against
\SI{18.7}{\percent} of currently deployed applications that use it (\S10). And the migration key's
scope --- whether \fluxd itself would reject a forced update touching anything beyond the version
field and the four renames, or whether that restraint would live only in the tooling that constructs
the message --- is open in \S10, as is whether the key is retired when the migration completes rather
than left live in a binary. Where the submitter is autonomous software rather than a person, the
specification would travel with a delegation certificate
$\mathsf{cert} = (pk_\agent, \scopeset, \budget, [t_0,t_1], \mathsf{rev}, \nu)$ signed by the
principal's full 2-of-2 ceremony, and a node would accept it by a predicate computable from the
message, its own chain state at $\hgt$, and nothing else (\S17). The deployment half of that
authority composes with machinery that already runs; the spending half does not, and \S17 classes it
as not achievable on the \Flux L1 without either a new trusted party or new consensus state.

\paragraph{2. The tenant would pay.}
One transaction, one recipient, one memo. The settled shape is a single transparent output paying the
full price $\price(\app,\hgt)$ to $\mathsf{sink}$ --- one application-payment address carried as a
consensus constant from $\hpa$ onward --- together with one zero-value \texttt{OP\_RETURN} whose data
field is the \num{32}-byte $H(m_\app)$; these are rules PS1 and PS2 of \S7
\srcintent{design intake} \srcdoc{reconciliation record}. There would be no
second output, no fee component and no new transaction type. The payer would construct one recipient
and its own change and would perform no other arithmetic, so any wallet able to pay an address and
attach a memo --- and any agent holding a bare wallet API --- could construct a valid payment without
linking a \Flux-specific library or waiting for tooling to be rolled out. \S7 records the deployed
\CumulusVPN entitlement path as the existence proof for exactly that shape, and records the two
shapes not taken and why: a fee-carried share would make the amount a function of the sending
wallet's change calculation, and a typed transaction would put a library dependency in front of every
payer. Direct per-deployment settlement would remain the default and the recommended path for
autonomous software, for the reason \S16 proves: entitlement is a function of the confirmed chain
prefix alone, so no account, no stored credential and no relationship outliving the deployment would
be involved. The alternative would be the custodial path \S16 documents as \shipped, in which a
facilitator holds the user's keys and pays the chain for them, and over which \S16 recommended that a
recurring charge be authorised by a \S17 certificate whose scope $\scopeset$ is restricted to
\emph{renew this application at no more than this price} under budget $\budget$, because that is the
only one of the three candidate authorities that removes the custodian rather than relocating it ---
a recommendation the settlement has since overtaken, since the protocol is to offer no automatic
renewal at all and recurrence stays in the service layer, by custody (\Cref{rem:no-protocol-renewal}).
Two of this step's open decisions were closed by the settlement: O3, the
memo's semantics, is $H(m_\app)$ in an \texttt{OP\_RETURN}, and O18, the serialisation and the
enforcement boundary, is the single-output shape with the split applied by consensus. Two remain.
O14, the identity of $\mathsf{sink}$, is unfixed --- keeping the address in use today preserves
tooling and one continuous history but mixes pre-$\hpa$ spendable receipts into the counter of
step~4, and a consensus-level address is a consensus-level commitment. O8, which payment channels
must settle through $\mathsf{sink}$ at all, is the load-bearing one, because every channel outside it
is revenue the consensus split never touches. Where a deposited balance would live --- on-chain,
\FluxOS-held or wallet-held --- was the architectural choice \S16 could not make, alongside the
authorisation object; both are now settled out of protocol scope, the balance held off-chain by a
third-party operator and no authorisation object existing at the protocol level at all
(\Cref{rem:balance-location}, \Cref{rem:no-protocol-renewal}).

\paragraph{3. The payment would be validated.}
Two consensus rules would run at $\hgt \ge \hpa$, and neither would involve the payer. PS3 makes any
output paying $\mathsf{sink}$ unspendable, so a block containing such a spend would be invalid and
the value would leave circulation at the moment of payment. PS4 divides what arrived:
$A_\Phi = \lfloor r(\hgt)\,v/100 \rfloor$ with $r(\hgt) = 100\,\nodeshare(\hgt) \in \mathbb{Z}$ is
credited to the fee pool, and $A_F = v - A_\Phi$ is credited to the Foundation, the single floor
taken once per received output and favouring the Foundation by strictly less than $1\sat$ (\S7).
\FluxOS would check three things and no others: that a payment exists, that it carries $H(m_\app)$,
and that it covers $\price(\app,\hgt)$. \textbf{The payer splits nothing and can therefore get
nothing wrong, and \FluxOS has no split to police.} That is the point \S7 makes about the decision
rather than about the outcome: choosing the single-output shape for \emph{wallet compatibility} also
removed the \emph{enforcement} ambiguity, and both properties came from one choice (C38). Three
consequences follow and \S7 proves them together. A wallet could not mis-round or omit a share that
it never computes. A \FluxOS release could not alter $\nodeshare(\hgt)$, the tier partition or the
destination of either share, because those would be evaluated in \texttt{ConnectBlock} rather than in
\FluxOS, and a node running a changed rule would be rejected by every other node. And the coinbase
would gain no output, because both shares would be added to outputs that already exist. What consensus
would \emph{not} fix is the amount: \FluxOS would still compute $\price$ from the height-scheduled
table, and \S7's formulation is the honest one --- \emph{consensus would fix the distribution, policy
would fix the base}, against a table revised five times with the CPU rate falling \num{100}-fold.
O7 --- whether received value can be visible to consensus while the payer is shielded --- is moot:
both shielded pools are to be retired (\S\ref{sec:shielded-retirement}), so all settlement is
transparent. Whether the table's rebasing becomes rule-based or oracle-driven is O9, settled as rule-based in
\S\ref{sec:pnr:recommended}, which leaves what the rule is open.
The signed negative acknowledgement of requirement R16.1, which would make an underpaid registration
observable to its payer, is specified in \S16 and is not built --- and it is the one repair this step
still owes an autonomous payer.

\paragraph{4. Value would enter the fee pool.}
This is the step with no deployed counterpart, and it is the mechanism the whole target economy rests
on. The fee pool would be a single consensus-state integer $\pool_\hgt \in \mathbb{Z}_{\ge 0}$ in
$\sat$ with $\pool_{\hpa-1} = 0$, \emph{committed in the coinbase} of every block from $\hpa$ and
recomputed by every validator from the parent value and the block body as
\[
  \pool_\hgt \;=\; \pool_{\hgt-1} \;-\; D_\hgt \;+\; \text{receipts}_\hgt ,
\]
with a block whose committed value differs from the recomputation being invalid (\S7, rule PC)
\srcintent{design intake} \srcdoc{reconciliation record}. That closes O2, and
the reasons \S7 gives for it are the ones that matter to a reimplementer rather than to a reader:
reorg-exactness would come free, because the pool would be a pure function of the chain, so a reorg
would recompute it with no rollback journal and no way for the pool to diverge from the UTXO set; no
new database would be introduced, because the coinbase parsing and enforcement path the commitment
needs is the one \fluxd runs every block for the dev-fund minimum and the three tier payouts
\shipped; and a divergent implementation would be rejected at the block where it first arises rather
than one block later. \S7 rejects a separate store on this network's own
precedent --- the \FluxNode registry needed a crash-recovery marker and a repair RPC precisely
because a side database can diverge from the coins database --- and rejects a header field as the
widest-blast-radius fork available.

Nothing would be added to the coinbase at any point in the mechanism. The Foundation's
\SI{80}{\percent} would flow through the \emph{existing} dev-fund output, the operator drip would be
added to the \emph{existing} three tier outputs, and the pool commitment itself would occupy the
coinbase's data-bearing field rather than a value-bearing one (\S7, I4). Supply would be conserved by
construction and not by bookkeeping: PS3 removes what is paid to $\mathsf{sink}$ from circulation and
the coinbase re-issues it, so circulating supply would differ from \S5's issuance schedule by exactly
the pool balance in flight, $\supply'(\hgt) = \supply(\hgt) - \pool_\hgt$, which at steady state is
\SI{30}{\day} of the operator share. The same rule turns the sink's balance into a
\emph{public cumulative-revenue counter}: the gross application revenue settled through this path
since $\hpa$ would be the balance of one address, queryable by anyone, publishable by no party in
particular and withholdable by none --- which is what would let \S\ref{sec:evaluation} replace an assumed $100\flux$
per application per month with a measurement, and what would make O17 closable by observation rather
than by disclosure. One qualification travels with the counter and must not be dropped: it would
measure only what settles through this path, so every channel left outside it by O8 is revenue the
counter would not see and the split would not touch.

Four things about the pool were undecided, all four are now settled (three at
\Cref{tab:pnr-recommended}, O21 at \Cref{rem:pnr:o21}), and two hard security requirements stand
regardless. O11, whether the pool is seeded at activation,
decided whether the mechanism ramps over roughly ninety days or switches on --- it is settled at no
seed, because a seed would be a one-off Foundation transfer to operators and would have to be
labelled one. O5 and O6, the destinations of a tier share with no payee and of the $\le 2\sat$
partition residue, are worth almost nothing per block and are nonetheless fixed explicitly --- both
stay in $\pool$ --- because an implementation difference in either is a chain split. O21, created by
the settlement itself, is settled on the same-block option (\Cref{rem:pnr:o21}): the division is
applied on arrival, at the $\nodeshare$ in force at the receipt height, and the Foundation share is
paid through the dev-fund output of the block that receives it, which needs no second accumulator
and imposes no lag but makes that output depend on the block body rather than on the parent state
alone --- a property the one-block deferral would have restored, and which is given up because it
serves a verifier \S\ref{sec:bridge} establishes cannot exist for \Flux. The requirements are S1 and
S2. An emergency block, produced through the hardcoded 2-of-4 key path with no time or stall gating,
must not drip, or the holders of two of four keys could accelerate the release of accumulated fees
--- $\pool_{\hgt-1}/\drip$ per block, paid to whichever nodes head the queue, including any they
operate --- on every block they produce; they could not redirect it, because an emergency block
still fills the deterministic tier payout (\S7). And forfeitures from the moderation mechanism of
step~11 must never enter $\pool$ (\S7).

\paragraph{5. Consensus would select the payees, and the pool would drip to them.}
The rotation would be unchanged: \PNR would consume \PoN's payee set and add no selection mechanism
of its own (\S7). Each block would drip $D_\hgt = \lfloor \pool_{\hgt-1}/\drip \rfloor$ from the
\emph{parent} state, partitioned by tier as
$D_{\hgt,t} = \lfloor 2\tbase{t}D_\hgt/27 \rfloor$ over the integers $2\tbase{t} \in \{2,7,18\}$,
giving $(\tierratio{\tC}, \tierratio{\tN}, \tierratio{\tS}) = (2/27,\ 7/27,\ 18/27) =
(\SI{7.41}{\percent},\ \SI{25.93}{\percent},\ \SI{66.67}{\percent})$, and the existing tier output
would carry $\subsidy_t(\hgt) + D_{\hgt,t}$ rather than a new output. The pool would update as
$\pool_\hgt = \pool_{\hgt-1} - D_\hgt + \epsilon_\hgt + U_\hgt + F_\hgt$. Two properties follow and
both matter for how a reader should read this step. At steady state the drip equals the inflow, so
the pool is a delay line with time constant $\drip$ blocks and not a treasury that anyone could
choose how to spend (\S7). And because the drip is taken from the parent state and is linear in it,
an operator who is also a tenant and times a payment to land the block before its own node is paid
recaptures exactly $\tierratio{t}/\drip$ of its contribution --- $7.72 \times 10^{-6}$ for a Stratus
node at the settled $\drip$, three orders of magnitude inside the design's own requirement (\S7). The
drip constant is settled at $\drip = \num{86400}$ blocks, exactly \SI{30}{\day} at \SI{30}{\second}
spacing, which closes O1 \srcintent{design intake}. The trade-off it resolves
does not disappear with the decision, and \S7 restates it so that a later revision starts from it:
within-tier spread $\approx \rot{\tC}/\drip$ and targeted recapture $\approx \tierratio{t}/\drip$
both fall with $\drip$, while lag, activation ramp and steady-state pool size $F\drip$ all rise with
it, over the domain $[10^{4},\ 2\times10^{5}]$ blocks in which the constant was chosen. Because
$\drip$ would be a consensus constant, revisiting it would be a network upgrade rather than a
configuration change --- so what accompanies the value is O12, settled as a documented review
trigger at $2\times$ the current Cumulus count rather than an automatic rule, since at $4\times$ the
decay-side spread reaches \SI{16}{\percent}.

\paragraph{6. The application would be placed across nodes.}
Placement would run under the \specv{9} reconciler: roughly \num{25} focused service packages
including \texttt{appLifecycle}, \texttt{appMesh} and \texttt{appPlacement}, a drain server exposing
\texttt{drain\_app}\slash\texttt{stop\_app} over a UNIX-socket JSON-RPC endpoint, and per-application-CA
backend TLS under which a node generates a local Ed25519 keypair and obtains a \SI{30}{\day} host
certificate signed by the application's own certificate authority through \fluxbench, renewed
automatically (\S10). That substrate is built and tested --- \num{289} unit-test files with
\num{7763} test cases run on every push --- and it lives on an individual contributor's fork rather
than in the organisation repository, which \S10 and \S23 both record as a fact about where
development authority sits (\S10). Two things are inert. \texttt{quorum\allowbreak{}Grant\allowbreak{}Activation\allowbreak{}Height}, the
height gate for the fork's mastership and active-standby quorum plane, is read by production code but
set only in integration-test fixtures and is absent from shipped configuration (\S10). And the
conflict-resolution rule itself would be unchanged: nothing in any drafted section replaces
earliest-broadcast-wins, so the target inherits both the \SI{90}{\second} wait and the honesty
assumption on \texttt{broadcastedAt} that \S9 states.

\proposednote{Earliest-broadcast-wins is retained at target. No section of this paper specifies a
replacement, and on reflection none is needed for the reason \Cref{rem:pnr:framing} gives: because
\PNR{} pays no marginal reward for hosting, winning a placement race carries no direct income, so the
race has no financial stakes and the incentive to game it is weak. A deterministic scheduler would be
preferable on other grounds --- predictability, capacity balancing --- but it is a \FluxOS{} design
question, not a consensus one, and replacing a mechanism that is not currently being gamed is not a
prerequisite for anything else in this paper. \srcinferred.}

\paragraph{7. The host would be attested.}
The proposed answer is not a better key but a design that stops depending on a key and depends on
collateral instead (\S14). A node would opt into the attested tier by posting a bond $\beta_i$
\emph{separate from} its \FluxNode collateral, which would remain self-custodied and untouched by any
slashing path. A claim would be a tuple
$\mathsf{cl} = (\mathsf{subject}, \mathsf{predicate}, \mathsf{value}, \hgt, \mathsf{nonce})$ signed under
the node's key; a committee of $n$ bonded nodes would be sampled deterministically as
$\mathcal{C}(\mathsf{cl}) = \mathrm{Top}_n(i \mapsto H(\mathsf{subject} \| \mathsf{prevHash}(\hgt) \| i))$,
reusing the sortition \fluxd already implements for \PoN slot eligibility; a claim would be accepted
at $k$ of $n$ identical signatures, challengeable for a window $\Delta_{\mathrm{chal}}$ by any party
posting $\beta_{\mathrm{chal}}$, and a resolved challenge would slash each original signer by a
fraction $s$ (\S14). Requirement~0 gates the entire construction: the attestation signing key would
have to be derived from or sealed to a per-machine hardware root, which
Property~``Attestation is not machine-bound'' in \S14 says it is not, and a tenant-reachable
verification path would have to exist where today there is none. Absent the first condition a bond
would secure a signature anyone holding the software can produce, making it forfeitable for
statements the bonded party never made. The four committee parameters are recommended by \S14 at
$n = 9$, $k = 6$, $\Delta_{\mathrm{chal}} = \num{2880}$ blocks and $s = 1$, and the bond-sizing
function at a floor $\beta_{\min}$ or a fixed fraction of the value secured, whichever is larger,
with both coefficients still to calibrate: $\beta_i$ must scale with the value secured by the claims
a node is sampled to sign, not with its \FluxNode tier, because a fixed bond cannot secure an
unbounded reserve (\S14). Admissible claim classes are likewise recommended as the objectively
re-checkable class only at launch, \S14 explicitly declining to invent a resolution procedure for
general computation-integrity claims.

\paragraph{8. \FDM and \PDNS would make the application reachable.}
This is the step where the target is least specified, and the paper should not pretend otherwise. The
requirement is legible from \S12's centralization inventory: \texttt{api.\allowbreak{}runonflux.\allowbreak{}io} is the sole
application-discovery source with no on-chain or peer-to-peer fallback coded anywhere, so any
decentralization of reachability starts with a discovery fallback that does not exist. Two further
interfaces would be needed and neither is designed. First, \S23 records as open whether \FDM and
\PDNS honour a ban by dropping routes at $\hgt_b + \Delta_{\mathrm{exec}}$, where $\hgt_b$ is the
height of the coinbase carrying the ban memo of step~11 --- without which a banned application would
remain reachable at its name until the last node had removed it, and with which the
Foundation-operated reachability plane acquires a new dependency on moderation state. Second, private overlay deployments (\VPON s) are a roadmap item with no wire
protocol, key-management scheme or placement-isolation mechanism specified anywhere; \S12 states the
two requirements a \VPON would have to meet --- membership gated by possession of a private key
rather than by resolvability, and a reachability plane not positioned to learn which nodes host a
private deployment or on whose behalf --- and states that the construction is open. The shipped
\CumulusVPN transports (AmneziaWG obfuscation, WireGuard-over-TLS, client-side multi-hop) are cited
by \S12 as building blocks a future design would not have to invent, not as a design.

\paragraph{9. It would run and be monitored.}
The target changes less here than anywhere else, and deliberately. Reservation pricing would remain,
so there would still be no metering on the application path, no counter to attest and no counter to
dispute; \S16 notes that this is the same choice \S7 makes one layer down when it removes any
metering or attestation dependency between payment and placement. What \specv{9} would add is
drain-aware shutdown: marking a component draining would trigger an immediate presence rebroadcast,
so the signal a routing layer consumes rides the gossip path \FluxOS already uses (\S10). The
companion \texttt{flux-shutdownd} pipeline is \inprogress --- its signal-and-cleanup stages are wired
to production events, while drain-broadcast and \texttt{preStop} execution are modelled and
unit-tested but not yet reachable (\S9). Whether the target architecture should acquire an
availability remedy at all is settled by \S16 as no, at the protocol layer (\Cref{rem:no-sla}): the
minimum viable input would be a third-party-observable availability record, which does not exist, and the
cost would be reintroducing the metering dependency the design removed.

\paragraph{10. It would renew or lapse.}
Two paths, and conflating them is the error this step exists to prevent. On the default path there
would be nothing to renew, because \textbf{the payment is the renewal}: a tenant wanting another term
would send another payment of the shape of step~2, and entitlement would remain what \S16 proves it
to be --- a function of the confirmed chain prefix alone. There would be no standing obligation, no
balance to track, no charge to authorise, and \emph{no lifecycle state machine at all}, so there
would be nothing for a policy to decide and no state a policy could get wrong. That is a consequence
of the design and not a gap in it: \S16's corollary is that a grace or suspension state would have to
be a function of information absent from the chain prefix, so the direct path cannot carry one
without surrendering the property that makes it the recommended path for autonomous software. Expiry
would stay arithmetic.

The second path would remain custodial, and \S16 documents what it is as \shipped rather than
projecting what it might become: a \emph{prepay-and-forward facilitator} that holds both of a user's
private keys --- the identity key whose address is written into the specification's \texttt{owner}
field, and the payment key --- charges a card, and pays the chain on the user's behalf, keeping no
user balance and no withdrawal path (C42). The four-state machine
$\mathrm{st}(\app,\hgt) \in \{\textsf{funded}, \textsf{grace}, \textsf{suspended}, \textsf{evicted}\}$
belongs to something else: a deposited-balance credits layer that \S16 specifies in full and marks
explicitly as unimplemented. Were it built, a renewal would fire at
$\hgt = e_\app - \Delta_{\mathrm{pre}}$ if the balance covers $\price^{\mathrm{ren}}$; a lapse would
move the application to \textsf{grace}, where it would keep running and keep its data; after
$G_{\mathrm{g}}$ blocks it would move to \textsf{suspended}, containers stopped but not removed and
volumes retained; and only after a further $G_{\mathrm{s}}$ blocks would it be \textsf{evicted} and
its data destroyed. The separation of stopping from deleting is that layer's design commitment and it
has a price \S16 bounds exactly: revenue forgone over the whole window is at most
$\ninst_\app\,p_{\mathrm{hdd}}(\hgt)\,\resvec^{\mathrm{hdd}}(\app)\,(G_{\mathrm{s}} +
G_{\mathrm{g}})/L^{\ast}(\hgt)$, linear in the declared disk footprint and independent of the CPU and
RAM reservation. Its four parameters are settled at \Cref{tab:credit-params} --- $G_{\mathrm{g}} =
\num{2880}$ blocks from the domain $[240, 2880]$, $G_{\mathrm{s}} = \num{20160}$ from
$[2880, 43200]$, $\Delta_{\mathrm{pre}} = \num{120}$ from $[20,\infty)$, bounded below by
confirmation depth plus propagation, and the resumption charge $\price^{\mathrm{res}}$ at one renewal
period from $[0,\infty)\flux$ --- and all four are service-layer values, because the two
architectural choices they were downstream of are settled out of protocol scope: a deposited balance
lives off-chain with a third-party operator, and nothing in the protocol authorises a recurring
charge (\Cref{rem:balance-location}, \Cref{rem:no-protocol-renewal}). One requirement crosses from
\S7 and constrains any answer: after $\hpa$ every
application-hosting charge, whatever the payment channel, must settle as exactly one L1 payment to
$\mathsf{sink}$, or the consensus split is void and the revenue counter of step~4 under-reports ---
which is why a credit balance may not become an off-chain ledger of hosting charges. And whichever
path a tenant took, the invariant that makes the custodial one optional rather than structural would
hold: for every deployment reachable through a custodian or through credits, the same deployment is
reachable at the same height through an explicit transaction the tenant signs and broadcasts itself
(\S16).

\paragraph{11. It could be reported and evicted.}
Every step of the decision would be a chain fact, and that is the whole difference from step~11 of
the first pass. A report would be a \fluxd transaction of a new \FluxNode-family type carrying the
application hash to be banned, shaped after the deployed start and confirm transactions and
\emph{indexed}, so that any node could enumerate the open reports on an application from its own
chain state; it would carry no fee and no bond, and its confirmation height $\hgt_0$ would fix the
canonical clock and \emph{freeze the electorate} (\S23)
\srcintent{design intake}. A window of
$\Delta_{\mathrm{vote}} = \num{2880}$ blocks --- \SI{24}{\hour} at \SI{30}{\second} spacing, chosen
short on the team's reasoning that removal often needs to be fast --- would open on that
confirmation. Each vote would be a second special transaction, signed with a node \emph{owner} key
and \textbf{fee-less}, which is what removes the disenfranchisement objection: no operator would need
spendable \flux at the key that carries its collateral in order to vote. \fluxd would count the
quorum internally from those transactions, so there would be \textbf{no off-chain tallier}, no
certificate for anyone to assemble and no gossiped object anywhere in the decision path --- the
property \S23 proves as public verifiability, under which two honest nodes holding the same canonical
chain hold the same enforcement set exactly rather than up to gossip delay. Weight would be
collateral-weighted at one unit per $\num{1000}\flux$, so per node $\tC \mapsto 1$,
$\tN \mapsto 12.5$, $\tS \mapsto 40$ against a measured total $\vtotal = \num{88514.5}$ units, and
the rule would be $Y_r(\hgt_c) \ge \vthresh \cdot \vtotal^{\!\star}(\hgt_0)$ with $\vthresh = 0.25$
as an \emph{absolute} fraction of total weight and abstention counting as ``no''.

On the threshold being crossed the block template would change: the \PoN producer would carry a
\emph{ban memo} in the coinbase, \textbf{at most one per block}, so bans would serialise at one per
\SI{30}{\second} slot and the coinbase would grow by at most one fixed-size memo however many reports
were open --- the same constraint \S7 respects when it routes the revenue split through outputs that
already exist. \FluxOS would observe the memo and purge, reusing the removal path the deployed
blocklist sweep already exercises, after a confirmation delay $\Delta_{\mathrm{exec}}$ that exists
because a purge is not locally reversible. A report that did not reach threshold inside its window
would simply expire: nothing removed, nobody penalised, the reporter's capacity released, and the
same application immediately re-reportable with no cooldown --- which \S23 argues is the mechanism
working rather than failing, because abstention counts against removal and most reports are expected
to end this way. Anti-spam would be priced in the resource that already prices voting: an identity
could hold at most as many concurrently open reports as it operates confirmed nodes, which is
Sybil-neutral because re-keying moves nodes between owner keys without changing the sum, caps the
network at $\lvert N(\hgt)\rvert = \num{6489}$ concurrent reports, and prices sustained nuisance at
$\num{1000}\flux$ of locked, recoverable collateral per slot --- with no bond, no custody script and
no forfeiture destination to argue about.

Two properties are worth carrying into this composition. Turnout is a liveness parameter and never a
safety parameter, because the threshold depends on the electorate and not on who votes, so low
turnout makes eviction strictly harder (\S23). And the safety property is exactly a collusion bound,
which at the measured distribution is $\kappa(0.25) = 4$ by owner identity --- the operative figure,
since votes are signed with owner keys --- falling to $3$ once addresses linked by a shared node key
or a shared collateral funding transaction are merged: safety against wrongful eviction would not be
a property of the mechanism but the assumption that three or four specific parties do not agree
(\S23). The category taxonomy was the single open parameter on which the mechanism's consistency with
the roadmap's own published ``behaviour, not content'' principle depended; its scope is now ruled on
--- the taxonomy is confined to behavioural harm, so the principle stands and the quorum is an
enforcement path for it (\Cref{rem:content-veto}) --- and only its enumeration remains open. The rest of what this composition first carried as
undecided is settled at the values of \Cref{tab:23-recommended}: the banned hash is over
$(\mathsf{app\_name}, \mathsf{owner})$, which is also the enforcement scope; inclusion of the ban
memo is consensus-required at every height at which a carried report remains unbanned, over the
first such report in the order least crossing height, then least $H(r)$, and a report on a hash
already in the ban set is invalid \srcintent{design intake};
$\Delta_{\mathrm{exec}} = \num{120}$ and $\Delta_W = \num{20160}$ blocks;
the relay policy for fee-less transactions is a per-identity mempool quota; the graduated
\SI{10}{\percent} and \SI{15}{\percent} tiers, each of which one or two identities can reach alone,
are removed; and $(m, T) = (21, \num{86400})$ for the recommended second chamber, which \S23 is
careful to say raises the cost of \emph{surprise} rather than the cost of \emph{capture}.
Reinstatement has no path at all, and \S23 states that as an open problem with requirements rather
than inventing one. Two questions could not be settled by choosing a parameter and had to be
verified in \fluxd, because the mechanism rests on them: whether the emergency-block path is subject
to the coinbase rules the ban memo would rely on --- if it were not, two of four hardcoded keys could
ban an application with no report, no votes and no threshold; it is, checked against source
(\Cref{rem:emergency-coinbase}) --- and whether re-addressing a node resets its recorded age, which
decides whether the tenure gate does any work and remains unverified. And one requirement crosses
back to \S7: forfeited tenant balances must never enter $\pool$, or
evicting an application would pay the fleet that voted to evict it.

The notation collision this composition exposed has been resolved: $\kappa$ now denotes \S23's collusion-set size and nothing else. \S14's claim object is written $\mathsf{cl}$ and the bridge's per-chain cap is $\overline{M}_c$ \srcdoc{reconciliation record}.

\paragraph{12. Operators would be compensated.}
From issuance \emph{and} revenue. At the proposed activation height $\hpa = \num{3466630}$ the
subsidy would be cut to $6.3\flux$ per block and the operator share would begin at
$\nodeshare = 0.20$, rising by $\rampstep = 0.05$ at each subsequent reduction to the cap
$\nodesharecap = 0.50$ at height \num{9378400} (\S7). At activation, operator issuance would be
$\num{6386040}\flux$ per year on the operator-only basis --- excluding the dev-fund remainder, which
is why \S7 corrects the design intake's own parity table downward by \SI{3.7}{\percent} --- and \PNR
income at the intake's activation-scale revenue assumption $V_0 = \num{1431600}\flux$ per year would
be $\num{286320}\flux$ per year, i.e.\ \SI{4.29}{\percent} of operator income. Parity revenue would
be $\num{31930200}\flux$ per year, $22.3\times V_0$. \S7 states the conclusion without softening:
\textbf{with flat demand there is no crossover, ever}, because the perpetual tail of
$\num{862659}\flux$ per year to operators exceeds $\nodesharecap V_0 = \num{715800}\flux$ per year
and the tail is perpetual by \S5. Two things were open at this step and both would have changed the
arithmetic; both are now settled. The sequencing of \PNR activation against the parallel-asset
emission cut, contested between the design intake and the public roadmap, is settled below, and O20
--- whether the PA-Depletion halving applies uniformly to the three tier bases and to the dev-fund
remainder --- is settled as uniform (\Cref{tab:pnr-recommended}), which is the basis every parity
figure above assumes. And the property that reorders the rest of this paper belongs here:
under the specified rule, a node's \PNR income would be a function of its tier and its queue position
and of nothing else, so its derivative with respect to applications hosted would be identically zero
(\S7).

Sequencing conflict C1a, which the transition table below depends on, is settled: \PNR{} activation
and the parallel-asset emission cut are \textbf{simultaneous} at $\hpa = \num{3466630}$
\srcintent{design intake} (\planned), superseding the roadmap's strictly-after
ordering (\Cref{rem:c1a-settled}). The accepted cost is stated with the decision rather than beneath
it: simultaneity means no external observer can confirm that \PNR{} pays operators before
parallel-asset mining rewards are retired, because there is no interval in which one has happened and
the other has not.

A numeric disagreement this composition exposed has been resolved in favour of \S7's basis: operator issuance at \PNR{} activation is $\num{6386040}\flux$ per year on the operator-only basis, and \S23's launch-deterrent proposition now uses it. The derived per-node deterrent of about $\num{0.10}\flux$ per Stratus node per application per year is unaffected at the precision stated \srcdoc{reconciliation record}.

\begin{remark}[The two retrieval heights are not a disagreement]
The measured chain tip appears as \num{2917115} in \S4, \S7's rotation remark, \S20 and \S23, and as
\num{2912015} in \S5's supply figure and \S7's activation-table baseline, both dated 2026-09-03.
C30 resolves this: the live API snapshot was taken at the first height and the emission model
evaluated at the second, roughly \num{5100} blocks earlier on the same date, so neither is wrong.
The convention this section follows, and which \S3's methodology paragraph states, is that every
figure names its own height rather than the paper fixing one.
\end{remark}

% =====================================================================
\subsection{What changes, and what does not}
\label{sec:e2e:diff}

%% longtable rather than table[H]: the twelve rows plus caption exceed one text page (162pt
%% overfull vbox, caption pushed off the page). Same header pattern as Appendix B. Round 18, G30.
\begingroup
\scriptsize
\renewcommand{\arraystretch}{1.2}
\begin{longtable}{@{}p{0.055\linewidth}p{0.265\linewidth}p{0.29\linewidth}p{0.30\linewidth}@{}}
\caption{The twelve steps, deployed against target, with the transition condition for each. The last
column is the one to read: since rounds 11--18 most rows terminate in a settled value awaiting code
or a network upgrade, and the rows that still end in a decision nobody has taken are 1 (the
migration key's scope and the envelope replacing \texttt{enterprise}), 2 (O14 and O8), 3 (R16.1,
specified and unbuilt, and the rebasing rule left open under O9) and 8 (a discovery fallback and a
\VPON{} construction).}
\label{tab:e2e-diff}\\
\toprule
step & deployed (\shipped) & target (\planned) & what has to be true for the transition \\
\midrule
\endfirsthead
\multicolumn{4}{l}{\small\itshape (Table~\ref{tab:e2e-diff} continued)}\\
\toprule
step & deployed (\shipped) & target (\planned) & what has to be true for the transition \\
\midrule
\endhead
\bottomrule
\endfoot
1 & \specv{8} enforced from height \num{1932380}; seven versions live; no \texttt{duration}, no
\texttt{paymentMemo} & \specv{9} with \texttt{duration}, \texttt{paymentMemo},
\texttt{referralCode}, \texttt{machineType}; \texttt{enterprise} removed; every application rewritten
forward by a one-off forced migration; agent submission under a \S17 certificate &
consumers written for the three economic fields; an enforcement height set \emph{after} them; an
encryption envelope replacing \texttt{enterprise}; the migration key's scope enforced in \fluxd
rather than by convention, and its retirement stated; a signature field and enforced canonicalization
in the schema \\
\addlinespace
2 & one transparent transaction to \texttt{t3Nryf\ldots}; hash in \texttt{OP\_RETURN}; price from a
six-segment table & \emph{same shape, moved into consensus}: one output of the full price to
$\mathsf{sink}$ plus one \texttt{OP\_RETURN} carrying $H(m_\app)$ (PS1--PS2); no second output, no new
transaction type & O14 ($\mathsf{sink}$'s identity) fixed; O8 (which channels must settle there)
fixed; O7 (shielded payers) moot after retirement; balance location and recurring-charge
authority settled as out of protocol scope \\
\addlinespace
3 & \FluxOS-only validation; underpayment silently dropped & PS3 makes $\mathsf{sink}$ unspendable
and PS4 applies the \SI{80}{\percent}/\SI{20}{\percent} split; \FluxOS checks only existence, memo and
coverage & a network upgrade; requirement R16.1 (signed \textsf{nack}) implemented; O9 (price-table
rebasing) decided \\
\addlinespace
4 & \textbf{does not exist} & $\pool_\hgt$ committed in the coinbase and recomputed by every
validator as $\pool_{\hgt-1} - D_\hgt + \text{receipts}_\hgt$; no new coinbase output; the sink
balance is a public revenue counter & O11 (no seed) and O5/O6 (both stay in $\pool$) settled; O21
(division on arrival, Foundation share paid in the same block) settled; S1 --- emergency blocks must
not drip; S2 --- forfeitures must not enter $\pool$ \\
\addlinespace
5 & \PoN sortition; one payee per tier per block; FIFO by last-paid height & unchanged rotation, plus
$D_\hgt = \lfloor\pool_{\hgt-1}/\drip\rfloor$ split by
$(\tierratio{\tC},\tierratio{\tN},\tierratio{\tS})$ & $\drip = \num{86400}$ is settled (O1 closed),
with O12 a documented review trigger at $2\times$ the Cumulus count and O10 drip ordering from the
parent state, both settled \\
\addlinespace
6 & opportunistic self-selection; broadcast, \SI{90}{\second}, earliest wins; Docker; reconciler &
\specv{9} reconciler, $\sim$\num{25} service packages, drain server, per-app-CA TLS &
\texttt{quorum\allowbreak{}Grant\allowbreak{}Activation\allowbreak{}Height} set in shipped configuration; earliest-broadcast-wins
retained (recommended above) \\
\addlinespace
7 & loopback mTLS only; key not machine-bound; benchmark not operator-bound; hash allow-list fails
open & bonded attestation network: bond $\beta_i$, committee $k$-of-$n$, challenge window
$\Delta_{\mathrm{chal}}$, slash $s$ & \textbf{Requirement 0}: hardware-sealed per-machine key and a
tenant-reachable verifier; $(n, k, \Delta_{\mathrm{chal}}, s) = (9, 6, \num{2880}, 1)$, $\beta_i$
scaling with value secured and re-checkable claims only are recommended, the bond coefficients still to
calibrate \\
\addlinespace
8 & \texttt{api.\allowbreak{}runonflux.\allowbreak{}io} sole source; 16-host \FDM fleet; HAProxy; Cloudflare or \PDNS &
requirements only: a discovery fallback; ban-honouring routes; \VPON membership by key &
a non-Foundation discovery source; a decision on \FDM/\PDNS honouring bans; a
\VPON construction, which is open \\
\addlinespace
9 & reconciler plus periodic monitors; reservation pricing; no metering & unchanged in kind; drain
signal over the existing presence path & \texttt{flux-shutdownd} drain-broadcast and
\texttt{preStop} reachable; no protocol availability remedy (settled, \S16) \\
\addlinespace
10 & expiry is arithmetic; no grace, no suspension, no data retention & \emph{unchanged on the
default path} --- the payment is the renewal, so no standing obligation and no state machine; funded
$\to$ grace $\to$ suspended $\to$ evicted specified for a credits layer and not deployed & for the
default path, nothing beyond step~2; for a credits layer, both architectural choices settled out of
protocol scope (balance off-chain with a third-party operator; no protocol renewal) and
$G_{\mathrm{g}}, G_{\mathrm{s}}, \Delta_{\mathrm{pre}}, \price^{\mathrm{res}}$ settled at
\Cref{tab:credit-params}; R16.2 honoured \\
\addlinespace
11 & unsigned GitHub file; $\kappa_{\mathrm{deployed}} = 1$; no record; \SIrange{7}{13}{\hour} to
fleet-wide effect & indexed report transaction; \SI{24}{\hour} window; fee-less votes; \fluxd counts
the quorum internally; ban memo in the coinbase, one per block; \FluxOS purges. $\vthresh = 0.25$
absolute, $\kappa(0.25) = 4$ & category enumeration fixed (its scope is settled as behavioural
harm); the remaining parameters settled at \Cref{tab:23-recommended}; both \fluxd{} verifications now
closed --- re-addressing does reset a node's age, and the emergency-path coinbase rules hold \\
\addlinespace
12 & issuance only; $\subsidypon$, tier bases $1{:}3.5{:}9$, dev fund $0.5\flux$ & issuance plus
revenue; $\nodeshare$ from $0.20$ to $\nodesharecap = 0.50$; \SI{4.29}{\percent} of operator income
at activation & demand growth, not price (parity at $22.3\times V_0$); C1a sequencing and O20
(uniform PA cut) both settled; and step~7 hardened first \\
\end{longtable}
\endgroup

\paragraph{The dependency that shapes everything.}
Read the table down the ``target'' column and the mechanisms look independent; read it down the last
column and one edge dominates all the others. \PNR would pay every confirmed node through the \PoN
rotation regardless of \emph{which} applications that node hosts, so the private stake any single
node has in aggregate network demand would be between $0.001\flux$ and $0.018\flux$ per year at
activation scale (\S7). That is not a free-riding opportunity --- placement is the network's choice
and an operator cannot decline it (\S7) --- but it does mean \PNR{} supplies no gradient rewarding
good hosting, and this is structural to the settled decision to pay all nodes rather than a tuning
error. What actually compels performance is not the payment but the \emph{eligibility gate}: a node
is a \PoN payee, and hence a \PNR payee, only if it is confirmed with a \fluxbench-signed transaction
and, at target, attested. \S4 establishes that the benchmark
signature is verified against five network-wide shared keys and is produced by a process on the
operator's own host, so it is not operator-bound; \S14 establishes that the attestation signature is
derived from constants compiled identically into every binary, so it is not machine-bound. The two
compose badly, and the composition is the reason this section exists.

\begin{property}[Revenue-funded compensation raises an existing exposure rather than creating one]
\label{prop:e2e-gate}
\planned. Under the settled \PNR constraints --- beneficiaries are all confirmed nodes, and there is
no metering or attestation dependency between payment and placement --- activating \PNR strictly
increases the expected payoff to occupying the paid set without providing capacity, by exactly the
amount \PNR pays: \SI{4.29}{\percent} of operator income at activation, rising as $\nodeshare$ ramps
toward $\nodesharecap$. The complementary bound is the one that decides the schedule. The other
\SI{95.71}{\percent} of that payoff is issuance already paid through the same gate today, so the
exposure is a property of \PoN{} and \PNR{} is a \SI{4.5}{\percent} increment on it. Hardening the
gate afterwards does not undo the interval; the interval it fails to undo is one twenty-third of the
standing one.
\end{property}

\begin{proof}[Proof sketch]
By \S7's zero-marginal-hosting property, a node's \PNR income is a function of its tier and its queue
position only, so the payoff to membership rises by the full drip while the payoff to hosting rises
by nothing. By \S4 and \S14 the predicate admitting a node to that queue is a benchmark signature the
operator's own host can produce under network-wide keys, together with an attestation that
demonstrates possession of a value every copy of the software contains. Hence the cost of satisfying
the predicate without providing the underlying capacity is unchanged by \PNR while the reward for
satisfying it rises, and the difference is the increase in expected payoff. The complementary bound
is arithmetic: operator issuance at $\hpa$ is \num{6386040}\flux\ per year against \PNR's
\num{286320}\flux\ (\S7), and both pass the same predicate. The step is a sketch and not a proof because it does not quantify
the cost side --- what it takes to occupy the paid set without providing capacity is not measured
anywhere in this paper, and \S7's Open Problem on hosting-conditioned eligibility is the construction
that would close it.
\end{proof}

The sequencing consequence is therefore not a preference, but it is also not the blocking dependency
an earlier draft of this paragraph asserted. Two things follow, and \S\ref{sec:migration} states the
schedule in these terms. Hardening the gate is the network's most urgent security work
\emph{independent of \PNR}, because almost all of the exposure it addresses is standing exposure that
\PNR{} did not create and will not remove. And \PNR's \emph{ramp} --- not its activation --- is what
should track that hardening, because the ramp is where the share of operator income sitting behind an
unbound gate actually grows. Placing the fix ``before or with activation, never after'' overstates
what the arithmetic supports and prices a \SI{4.3}{\percent} supplement as though it were the whole
exposure.

\paragraph{What the reader should notice about the two passes.}
Read as one document, the two passes make a single claim, and it is the thesis this paper has been
earning for twenty-three sections. \Flux today is a working decentralized cloud --- several thousand
independently operated machines, in dozens of countries, running \num{1195} tenants' containers, paid
for on-chain, with a consensus layer that produces a block every thirty seconds without a miner ---
whose \emph{trust}, as opposed to whose compute, is concentrated in operated infrastructure and
shared keys: one discovery API with no fallback, sixteen Foundation-run reachability hosts, one
Foundation-private network segment, one GitHub file as the moderation root, five network-wide
benchmarking keys, one release-signing key inside one HSM, a hardcoded 2-of-4 key set that can
produce blocks with no gating, and one transparent address that receives every tenant payment on the
network. The target architecture does not add capacity, features or throughput to that picture in any
material way; what it does, step by step, is move each of those trust concentrations into one of
three things the protocol can enforce --- \emph{collateral} (the quorum's weight, the report-capacity
bound, the attestation bond), \emph{attestation} (a machine-bound key and a verifier a tenant can
reach), and \emph{consensus rules} (the split as a rule instead of a promise, the pool as consensus
state committed in the coinbase, the eviction as a tally every validator recomputes from block data
and a memo in a coinbase rather than a document somebody publishes). That is the whole architectural
bet, stated once: not that the cloud gets bigger, but that the trust moves from parties into rules.
Every open parameter in the last column of Table~\ref{tab:e2e-diff} is a place where that move is
specified but not yet made.

\paragraph{The sequence the second pass can state and the first cannot.}
\planned. One consequence of the settlement decision is concrete enough to write as a sequence, and
it is the agent-native argument of \S18 reduced to three steps with nothing left implicit. A tenant
holding \emph{any} wallet and \emph{no account} would complete steps~1 to~3 unaided: construct a
\specv{9} specification and sign it with a key it generated itself; broadcast one ordinary
transparent transaction paying one address, with one \texttt{OP\_RETURN} attached; and have consensus
divide the proceeds, at which point every node independently agrees the deployment is funded to a
particular height. No second recipient to compute, no split to get right, no Flux-specific library to
link, no credential to be issued, no custodian, no relationship that outlives the deployment. Each
step's dependency is named where it falls rather than accumulated here --- step~1 needs consumers for
the three economic fields and an enforcement height after them, step~2 needs $\mathsf{sink}$ fixed by
O14, and step~3 needs R16.1 so that a short payment produces a signed answer rather than silence,
which for software with nobody to escalate to is the difference between a retry and an unrecoverable
loss. The first pass cannot state the same sequence, and the reason is recorded at its own step~3.

\paragraph{What would falsify it.}
A thesis that cannot fail is not research, so the specific, datable observations that would show the
transition is not working are named here rather than left implicit. \emph{First}, \PNR remaining a
rounding error: application revenue staying near the intake's $V_0 = \num{1431600}\flux$ per year, so
that \PNR stays close to \SI{4.29}{\percent} of operator income and the $g = \SI{0}{\percent}$ row of
\S7's crossover table holds --- under which there is no crossover ever, and revenue-funded
compensation remains a supplement rather than the mechanism. The settlement decision is what makes
this the cheapest item on the list to check: after activation the quantity is the balance of one
consensus-known, consensus-unspendable address, read from any block explorer, rather than a figure
anyone has to publish. \emph{Second}, the attestation key staying non-machine-bound:
\texttt{/v2/attest} continuing to derive its signing key from values compiled identically into every
binary at the point where bonds are first posted, which would mean Requirement~0 was declared met
without being met and that a bond secures a signature anyone with the software can produce
(\S14). \emph{Third}, \specv{9} never receiving an enforcement height: the public pricing endpoint
continuing to return no key for version~9 while \texttt{paymentMemo}, \texttt{referralCode} and
\texttt{machineType} keep zero consumers, and no forced-migration message appearing on any owner's
application --- which would strand \PNR, credits, referrals, the VPS tier and the provisioning API
together, since \S\ref{sec:feasibility} places all five on the same critical path. \emph{Fourth}, the shielded pools not
being retired: \texttt{Contextual\allowbreak{}Check\allowbreak{}Transaction} still rejecting transparent-to-shielded
transactions with \texttt{bad-\allowbreak{}txns-\allowbreak{}fluxrebrand-\allowbreak{}vprivacy-\allowbreak{}not-\allowbreak{}empty}, pool totals still at or near the
measured $\num{96479.94}\flux$, and no sunset height announced, at any date past
$\hshield$ (\S\ref{sec:shielded-retirement}). Note that this test was inverted by the
retirement decision: an earlier draft made \emph{failure to reopen} the falsifying observation, and
the observable that would have confirmed the old commitment --- a transparent-to-shielded transaction
confirming on mainnet --- would now falsify the new one. The paper records the inversion rather than
quietly re-aiming the test, because a prediction that changes direction is exactly where a reader
should check that the paper is tracking decisions rather than rationalising them.
\emph{Fifth}, a moderation capture event: a ban
memo in a coinbase whose crossing tally is carried by owner identities that are a subset of the four
largest --- three under the strongest linkage computable from public data --- or, equally decisive
and easier to observe, an application disappearing from the fleet with no report transaction, no
votes and no ban memo behind it, at any date after the quorum is said to be live. Each of these is
checkable from public data by a third party, on a stated date,
without access to any repository. That is the standard this section holds itself to, and it is the
standard by which the rest of the paper should be judged.

%% Drain this section's float queue before the next begins.
\FloatBarrier
